\section{Discussion}

\paragraph{Faithfulness as a family of constructs.} This study set out to determine whether faithfulness measurements are stable across classification methodologies. The central finding is that they are not: three classifiers applied to the same 10,276 influenced reasoning traces produce overall faithfulness rates spanning a 12.9-percentage-point range, per-hint-type gaps as large as 43.4 percentage points, and model rankings that reverse depending on the classifier used. CoT faithfulness, as currently operationalized in the literature, does not appear to be a single measurable quantity but rather a family of related constructs whose values depend critically on the measurement instrument. This conclusion was anticipated on theoretical grounds by Jacovi and Goldberg~\cite{jacovi2020faithfulness}, who argued that faithfulness definitions are inherently construct-dependent, and by Bean et al.~\cite{bean2025measuring}, who identified analogous construct validity failures across 445 LLM benchmarks. The present results provide direct empirical confirmation at scale.

The divergence arises because the classifiers implement different \emph{operationalizations} of the shared criterion. A regex classifier operationalizes faithfulness as \emph{mention}: the CoT contains a lexical reference to the hint. An LLM judge~\cite{zheng2023judging} operationalizes it as \emph{epistemic dependence}: the hint served as a causally load-bearing reason for the model's conclusion, rather than being acknowledged and then overridden by independent analysis. These are distinct constructs that happen to share the same label. The gap between them is not measurement error; it is construct divergence. The hint-type pattern reinforces this interpretation: sycophancy hints exist on a continuum between social acknowledgment and epistemic dependence~\cite{sharma2023sycophancy}, creating maximum ambiguity between the two constructs and producing the largest classifier gaps. Hu et al.~\cite{hu2025monica} similarly found that sycophancy in reasoning models resists binary classification and requires real-time calibration. Grader hints, by contrast, tend to be either fully embraced or fully ignored, leaving little room for interpretive disagreement. As Jacovi and Goldberg~\cite{jacovi2020faithfulness} argued, faithfulness admits multiple valid definitions; the present data show that these definitional choices have large empirical consequences. Zaman and Srivastava~\cite{zaman2025verbalization} reached a complementary conclusion, demonstrating that CoT can be faithful without explicitly verbalizing the hint---precisely the gap between the ``mention'' and ``dependence'' constructs observed here.

This pattern has precedent outside faithfulness evaluation. Gordon et al.~\cite{gordon2021disagreement} demonstrated that correcting for annotator disagreement in toxicity classification reduced apparent model performance from 0.95 to 0.73 ROC AUC, a degradation comparable in magnitude to the gaps observed here. Plank~\cite{plank2022problem} argued more broadly that such label variation reflects genuine interpretive plurality rather than noise, and should be modeled rather than suppressed. The present study extends these insights from human annotation to automated classification of reasoning traces, suggesting that a similar interpretive ambiguity that produces inter-annotator disagreement in subjective tasks also produces inter-classifier disagreement in faithfulness evaluation.

The construct multiplicity is further evidenced by the diversity of definitions in the faithfulness literature itself. Turpin et al.~\cite{turpin2023unfaithful} defined unfaithfulness as producing ``wrong reasoning'' that omits known influences; Chen et al.~\cite{chen2025reasoning} defined it as failure to ``acknowledge'' a hint; Lanham et al.~\cite{lanham2023measuring} operationalized it via causal interventions on intermediate steps; Lyu et al.~\cite{lyu2023faithful} and Radhakrishnan et al.~\cite{radhakrishnan2023decomposition} proposed structural decompositions to improve it; and Shen et al.~\cite{shen2025faithcotbench} benchmarked instance-level faithfulness across tasks. Each captures a different facet of the same underlying concern, and each would yield different numbers on the same data. Ye et al.~\cite{ye2026faithfulness} provide complementary mechanistic evidence that faithfulness decays over long reasoning chains, suggesting that classifier sensitivity may interact with trace length in ways not yet characterized.

\paragraph{Practical recommendations.} Four practices are proposed for future CoT faithfulness evaluations:

\begin{enumerate}
    \item \textbf{Report classifier methodology.} Every faithfulness number should be accompanied by a precise description of the classification approach, including the full prompt for LLM-based judges. This parallels broader best practices for subjective labeling tasks, where label definitions and annotation procedures must be specified explicitly~\cite{plank2022problem, gordon2021disagreement}.
    \item \textbf{Report sensitivity ranges.} Rather than single point estimates, report bounds: e.g., ``faithfulness is 69.7\%--82.6\% depending on classifier stringency.'' This communicates the inherent measurement uncertainty.
    \item \textbf{Cross-classifier agreement.} Applying at least two classifiers to the same data should become standard practice. High agreement increases confidence; low agreement signals that the result is classifier-dependent.
    \item \textbf{Caution in cross-paper comparison.} Faithfulness numbers from papers using different classifiers should not be treated as directly comparable without sensitivity analysis or careful adjustment for methodological differences. DeepSeek-R1's 39\%~\cite{chen2025reasoning} and the present study's 74.8\% (Sonnet) and 94.8\% (pipeline) are not contradictory findings; they are measurements of different constructs.
\end{enumerate}

\paragraph{Cost-accuracy trade-offs.} The three classifiers span a wide cost range: regex-only is free and instantaneous but limited to surface-level lexical matching (74.4\% overall); the regex + Ollama pipeline incurs zero marginal cost on a local GPU with moderate coverage (82.6\%); and Claude Sonnet~4, the most expensive option at \$48.99 for 10,276 cases (\textasciitilde\$0.005/case), produces the strictest judgments under this criterion (69.7\%). Counterintuitively, the cheapest and most expensive classifiers agree more closely with each other (74.4\% vs.\ 69.7\%, a gap of 4.7 percentage points) than either agrees with the pipeline (82.6\%). This pattern arises because the pipeline's local LLM judges surface cases that regex misses but that Sonnet deems non-load-bearing, inflating the pipeline's acknowledgment rate relative to both endpoints. For practitioners operating under budget constraints, regex-only classification provides a low-cost, surface-form baseline. When resources permit, running two classifiers of differing stringency and reporting the resulting range is preferable to relying on any single number. The \$48.99 total cost of the Sonnet judge is negligible relative to typical inference budgets, making dual-classifier evaluation accessible for most research groups.

\paragraph{Limitations.} This study compares three automated classifiers but lacks human ground-truth labels, which would provide an additional anchor for evaluating classifier validity. The Sonnet judge may also exhibit systematic biases identified in the LLM-as-judge literature~\cite{gu2024survey, ye2024justice}, though the binary classification task used here is less susceptible to position and verbosity biases than pairwise comparison settings. Human annotation was infeasible at the scale of 10,276 cases but would be valuable for a targeted subsample in future work. More broadly, this study compares only three predominantly text-based classifiers; other approaches, including human annotation with adjudicated disagreement analysis, causal interventions (e.g., removing hint-related CoT segments and observing answer changes), representation-based methods (probing or attribution on internal activations), and calibration-style scoring (degree of dependence rather than binary labels), may yield different sensitivity patterns. The observed disagreement may therefore understate the full range of classifier sensitivity across the broader space of faithfulness measurement methods. Additionally, all three classifiers operate on the visible CoT; they cannot assess whether the model's internal reasoning process actually relied on the hint, only whether the text acknowledges it, a fundamental limitation shared by all text-based faithfulness measures~\cite{arcuschin2025wild}. Recent work on steganography~\cite{roger2023steganography} and alignment faking~\cite{greenblatt2024alignment} suggests that models may encode reasoning in ways not captured by surface-level text analysis, and Marks et al.~\cite{marks2025auditing} demonstrated that hidden objectives can persist even under careful monitoring, further underscoring the need for classifier-robust evaluation methods.

\paragraph{Connection to thinking-token monitoring.} These results motivate approaches that are less sensitive to classification methodology. Baker et al.~\cite{baker2025cot} argued that CoT may be highly informative despite being unfaithful at the surface level, Yang et al.~\cite{yang2025monitorability} investigated the monitorability of CoT in reasoning models, and Xiong et al.~\cite{xiong2025thinking} measured faithfulness specifically in thinking drafts. A companion study of thinking-answer divergence~\cite{young2026suppression} finds that the gap between a model's internal thinking tokens and its visible CoT may provide a more robust signal of unfaithful reasoning, because thinking-token analysis does not require subjective judgments about whether a hint reference is ``load-bearing.'' Classifier-independent signals may ultimately prove more reliable than any single faithfulness metric~\cite{zhang2025probing}.
